\section{Estimation and inference}
\label{sec:inference}

\subsection{The estimator, and uncertainty that includes the model}
\label{sec:boot}

The pipeline is frozen and identical on every log, and the \emph{pipeline}
axis spans encodings as well as model families; every learner, encoder and
calibrator is the \texttt{scikit-learn} implementation
\citep{pedregosa2011scikit} at the version \texttt{requirements.lock} pins.
Attributes are assigned to roles by the registered rules of
Section~\ref{sec:design}, cases are ordered by their first timestamp, and the
split is temporal (Supplement~\ref{app:pipeline}). The usual interval for
such an increment resamples test rows with the models held fixed, which
excludes training-sample, model-selection, encoder and split variability and
the temporal dependence between cases. Every interval here is instead a
\emph{nested moving-block bootstrap}: one draw resamples the training half by
moving blocks \citep{kunsch1989jackknife} of length $\lceil n^{1/3} \rceil$,
refits every arm of every rung under every learner and quality mechanism,
and evaluates on a moving-block resample of the test half, with encoders,
frequency tables, target encodings, register rankings and calibrators
rebuilt inside the draw. On the simulation of Section~\ref{sec:sim} the
fixed-model interval covers a median \coverageNaiveMedian\ and the nested
basic interval \coverageNestedBasicMedian\ --- about \nestedGainInSe\ Monte
Carlo standard errors apart, which that design cannot resolve --- so the
nested construction is retained because the variability it admits is
variability the estimand contains, not because it measures better here; the
price it charges is the subject of the next two paragraphs.

\noindent Throughout, $K$ denotes the \textbf{cardinality of the register
field} --- the number of distinct values $f$ takes on a pair, \cardFMedian\ at
the median pair and \cardFMax\ at the largest (Table~\ref{tab:roles}) --- and
$n$ the number of training cases. The ratio $K/n$ governs everything in this
subsection and in Section~\ref{sec:regions}.

\paragraph{Refitting inside the draw breaks the percentile interval}
A bootstrap resample holds about $1 - e^{-1}$ of a high-cardinality
categorical's distinct levels, so every refit sees a sparser register and
returns a smaller increment; the bootstrap distribution is displaced and a
percentile interval taken from it is a correct interval for the wrong
quantity. On the sparse world of Section~\ref{sec:sim} it covers
\coverageSparsePct\ against a nominal $95\%$, \emph{worse} than the
fixed-model interval refitting was meant to improve on. Every interval here
is therefore the \emph{basic} (pivotal) interval
$[\,2\hat{V} - q_{1-\alpha/2},\; 2\hat{V} - q_{\alpha/2}\,]$
\citep{davison1997bootstrap}, and simultaneous bands are recentred by the
same shift (Supplement~\ref{app:pctfail}).

\paragraph{Weights rather than subsets, and what the comparison found}
Every band this paper reports comes from a bootstrap that draws Dirichlet
\emph{weights} per moving block rather than a subset of rows, so every row
has positive weight in every draw and every register level is present in
every refit. The multinomial alternative was run beside it over the same
design, at the same draw count and with the same seeds, on \schemeCells\
cells, and Supplement~\ref{app:schemes} reports the comparison in full; two
findings carry into the article. Under a multinomial resample a draw can drop
an arm entirely, so the cell gets no band at all and is counted as
\emph{unresolved} beside cells that were banded and straddled zero ---
\nUnbandableSameDesign\ of \nUnbandableSameDesignOf\ cells on this design,
against \nUnbandableCells\ under weights. And level loss is not what
displaces the bootstrap distribution: weights raise the share of the
register's levels a draw contains from \levelShareMult\ to
\levelShareWeighted, and the displacement's median moves only from
\dispMedianMult\ to \dispMedianWeighted, while the share of cells whose
pivotal interval excludes its own estimate falls from
\shareExcludesMultPct\ to \shareExcludesWeightedPct\ and not to zero. The
displacement that remains is unexplained: the level-loss mechanism of
Supplement~\ref{app:pctfail} is the reason the percentile interval is not
used, and it is not a diagnosis of the residual, which
Section~\ref{sec:lim} carries as an open defect.

\subsection{Three surfaces, and which one a claim may quantify over}
\label{sec:threesurfaces}

A statement of the form \emph{this surface is uniformly beneficial} is a
statement about a denominator, and the denominator has to be the set the
statement ranges over. Three sets are distinguished, named, and counted from
the files rather than quoted (Table~\ref{tab:denominator}). The
\textbf{declared surface} is every scientifically admissible specification:
the product of the axis levels declared in Section~\ref{sec:design}. The
\textbf{computational surface} is every specification actually evaluated; here
it equals the declared surface --- \nComputationalCells\ scalar cells against
\nDeclaredAdmissibleScalar\ declared, with \nDuplicateCells\ duplicates and
\nMissingCells\ missing --- and the audit script asserts that identity rather
than the manuscript asserting it. The \textbf{inference surface} is every
specification carrying bootstrap draws; it is smaller for one reason, that a
draw refits the whole pipeline and its cost is the declared surface's
multiplied by the draw count.

\textbf{The rule that selects it reads neither a result nor the log}: one
design on every pair --- two learners, two splits, three register-quality
conditions and three admissible rungs, crossed, \nCellsPerDecomposition\
cells a pair and \nInferenceScalar\ scalar family members across the corpus,
at \nDrawsSurface\ draws --- and the audit checks rather than declares that
it is a complete factorial, each combination present exactly once on every
one of the \nPairs\ pairs, so a corpus median over these families is a
median over one design (Supplement~\ref{app:pipeline}).

\textbf{A design, not a corner.} Because the inference surface is a balanced
factorial, the decomposition of Section~\ref{sec:sobol} can be computed on it
as well as on the declared surface --- one decomposition per pair per
instrument, over \nCellsPerDecomposition\ cells each --- and
Supplement~\ref{app:infanova} reports that the axis carrying the largest
first-order index agrees between the two surfaces on a minority of pairs.
That is a comparison between two designs that declare different levels, and
so a statement about the levels an analyst declared rather than a verdict on
either surface.

\noindent What the design does not change is the thing that matters most
for reading a region label. The inference family is a fraction of
the admissible surface --- a median \inferenceShareAdmissibleMedianPct\ of a
pair's admissible cells, and less than a twentieth on the pair with the
richest grid --- and \textbf{a label computed over fewer cells is
systematically cleaner than one computed over all of them}.
Section~\ref{sec:regionsdef} says which label that most flatters, and
Section~\ref{sec:regions} reports where it lands.

Every quantity that does \emph{not} need a bootstrap --- the decomposition,
the sign-disagreement rate, the regret comparison --- is computed on the full
declared surface; only the region labels are computed on the inference
surface.

\subsection{Simultaneous bands over the family the claim ranges over}
\label{sec:simbands}

A decision curve read across \nGrid\ thresholds, a ladder read across rungs,
or a surface read across every admissible cell, is a family of statements, and
a pointwise interval does not control the probability that \emph{some} member
of the family is wrong. For a declared family $\mathcal{F}$ of cells with
bootstrap draws $V^{(b)}_c$ and per-cell standard errors $\hat\sigma_c$, let
\[
T^{(b)} \;=\; \max_{c \in \mathcal{F}}
\frac{\bigl|V^{(b)}_c - \bar{V}_c\bigr|}{\hat\sigma_c},
\qquad
q_{1-\alpha} = \text{the } (1-\alpha)\text{ quantile of } \{T^{(b)}\},
\]
and report $\hat{V}_c \pm q_{1-\alpha}\hat\sigma_c$. This is the max-$t$
(Romano--Wolf) construction \citep{romano2005stepdown} and has family-wise
coverage over $\mathcal{F}$ asymptotically, and over nothing larger.

\paragraph{The family must be the family}
Within-cell bands, each covering five members, do not cover the several
hundred cells that a surface-level label ranges over. Four interval objects are
declared and every table says which it prints; \textsc{whole-surface} ---
max-$t$ over \emph{every} scalar cell of the pair's inference family, a
median \inferenceShareAdmissibleMedianPct\ of its admissible scalar cells
(Section~\ref{sec:threesurfaces}) --- is the only family a surface-level
claim may use, and it has a median
\familySizeMedian\ members against the within-instrument family's five, at a
median critical value of \maxTSurfaceMedian\ against that family's
\maxTWithinMedian\ and the pointwise $1.96$.
Family control is per log--target pair; every cross-pair statement here is
descriptive.
Supplement~\ref{app:families} gives these four \emph{interval} objects
--- which are not the reporting objects of Table~\ref{tab:objects} --- and the
rule for dropping a replicate that leaves a cell without outcome variation.

\paragraph{Estimating the critical value}
Two estimators of $q_{1-\alpha}$ are computed on every family. The first is
the empirical $(1-\alpha)$ quantile of the $B$ observed maxima $T^{(b)}$ ---
the Romano--Wolf construction as cited --- whose precision is that of an
order statistic of $B$ values: its order-statistic interval has a median
width of \qEmpWidth\ on this corpus's families. The second standardises the
draws and multiplies them by independent standard normals, which gives a
statistic with their empirical covariance that can be generated
\nMultiplier\ times for the cost of matrix arithmetic
\citep{chernozhukov2013gaussian}, so the budget fixes $B$ and not the
precision of $q$ (a Monte Carlo interval of width \qMultWidth). The
multiplier is an approximation: it reproduces the covariance of the draws and
therefore only the excursions a Gaussian process makes, and on every family
here it lies below the empirical quantile's whole order-statistic interval
(Supplement~\ref{app:twoq}). \textbf{The operative critical value is the
empirical quantile}, at a median \maxTSurfaceMedian\ on the whole-surface
families against the multiplier's \maxTSurfaceMultMedian;
Section~\ref{sec:simband} is the measurement that decides it, and every
table prints what the multiplier would have resolved beside what the
reported band does. A max-$t$ band also inherits the coverage of the
interval it is built from, which Section~\ref{sec:simsensitivity} shows is
governed by the ratio of register cardinality to training size;
Supplement~\ref{app:calibration} derives a widening calibrated at each pair's
own ratio and reports it as a sensitivity.

\paragraph{What these bands cover, measured against a known answer}
Section~\ref{sec:simband} measures family-wise coverage over synthetic
families matched to this corpus in size, draw count, cross-cell correlation
and tails, whose truth is known by construction and whose bootstrap is
\emph{ideal}, so every coverage it reports is an upper bound on what the real
procedure attains. On the family matched to the design this paper reports
--- \nCellsPerPairMedian\ cells at \nDrawsSurface\ draws under this corpus's
tails --- the empirical quantile's band covers \covEmpHeavyWholeAtDesign\
against a nominal $95\%$, at a Monte Carlo standard error of
\covEmpHeavyWholeAtDesignSE, and \covEmpNonzeroWholeAtDesign\ under a
non-zero truth, the regime in which a cell can be resolved correctly and a
region label is actually being asked to hold. The multiplier's band covers
\covMultHeavyWholeAtDesign\ and \covMultNonzeroWholeAtDesign\ on the same
families, several standard errors short, and no draw budget closes that gap:
along a ladder of draw counts its coverage plateaus from \nBootOneFifty\
draws on, while the empirical quantile's rises to its level by
\nDrawsSurface. \textbf{The level is bought with width}: the empirical band
is \widthEmpOverMultAtDesign\ times as wide on the matched family and
\ratioCorpusWhole\ times on the corpus's own families, and it resolves
\nResolvedWhole\ cells corpus-wide where the multiplier would resolve
\nResolvedWholeMult\ (Section~\ref{sec:regions}). What the measurement does
not cover is stated with it: the bootstrap is ideal there, so the
$K/n$-governed shortfall of the pointwise interval underneath
(Section~\ref{sec:simsensitivity}) and whatever the real resampling adds
are outside it, and Section~\ref{sec:lim} carries both.

\subsection{Planned contrasts, and a $p$-value the design can produce}
\label{sec:planned}

Five contrasts on the primary log were fixed before this analysis round ---
that the register adds to the intake baseline, that it still adds once the
free field is admitted, that it adds at the later decision time, that the
absorption $D$ is positive, and that a half-empty register is worth less than
a full one --- each with its complete estimand in Table~\ref{tab:confirm}.
They are not confirmatory, because they were fixed after the data were seen,
so their Holm-adjusted one-sided bootstrap $p$-values \citep{holm1979simple}
are reported for legibility: \nPlannedReject\ reject at $\alpha = 0.05$, but
a contrast is called resolved only when its interval says so, and on
\textsf{PC3} the one-sided $p$-value and the two-sided basic interval
disagree, so the count under this paper's own rule is \nPlannedResolved\
(Supplement~\ref{app:pcdetail}, which also gives the smallest attainable
adjusted value, \pMinAttainableHolm\ at \nDrawsPlanned\ draws). Every other
comparison in this paper is \emph{exploratory} and carries no multiplicity
correction, because none would be meaningful over a surface enumerated in
full.

\subsection{The sensitivity decomposition}
\label{sec:sobol}

Because the surface is a complete factorial over its axes --- the audit of
Section~\ref{sec:threesurfaces} establishes that it is, cell for cell --- the
variance of $V_s$ admits an exact and \emph{orthogonal} functional-ANOVA
decomposition \citep{sobol2001global,saltelli2010variance} under a product
measure $\mathbb{P} = \bigotimes_i w_i$, with first-order and total indices
$S_i$ and $S_{T_i}$ computed exactly rather than estimated
(Supplement~\ref{app:sobolest}). The closest
precedent is Hutter et al.~\cite{hutter2014efficient}, who decompose an algorithm's
performance over its \emph{configuration} space by the same device; here the
axes are analysis decisions, the components are exact over a complete
factorial rather than estimated from a surrogate, and the decomposition is
one of three reporting objects rather than the conclusion.

\paragraph{The split is an error stratum, and the primary decomposition treats
it as one}
Five of the split axis's six levels are expanding-origin folds on the same
data. They are different \emph{samples}, not different analyses, and a
decomposition that pools them reports sampling variability as specification
sensitivity. Because the decomposition is exact and orthogonal, the repair is
a partition rather than a model: every component either involves the split or
does not, so
\begin{equation}
\label{eq:strata}
1 \;=\; \underbrace{\textstyle\sum_{\emptyset \neq u \subseteq \mathcal{A}} S_u}_{\text{analyst choice}}
\;+\; \underbrace{S_{\{\text{quality}\}}}_{\text{counterfactual}}
\;+\; \underbrace{\textstyle\sum_{u \not\ni \text{split}} S_u - (\cdots)}_{\text{mixed}}
\;+\; \underbrace{\textstyle\sum_{u \ni \text{split}} S_u}_{\text{resampling}},
\end{equation}
where $\mathcal{A} = \{\text{pipeline}, \text{rung}\}$ and the third term
is what the split-free components carry beyond the first two,
and the four terms are computed, not estimated (Table~\ref{tab:strata}); the
identity holds to \stratumPartitionError. The first three are what survives
fold-averaging, and they are decomposed again \emph{on} the fold-averaged
surface, where each index is a share of the variance that fold-averaging
leaves and the higher-order share is net of resampling.
\textbf{Section~\ref{sec:axes2} reports that decomposition as the primary one}
and the pooled decomposition beside it, because the pooled one is what a
reader of the earlier specification-curve literature will expect and the
difference between them is itself a result.

The per-axis difference $S_{T_i} - S_i$ is the variance axis $i$ is
\emph{involved} in beyond its main effect. These overlap --- a two-way
component appears in both its axes' total effects --- so their largest member
is not the share carried by interactions. That share is
\begin{equation}
\label{eq:inttotal}
S_{\mathrm{int}} \;=\; 1 - \sum_i S_i ,
\end{equation}
which this paper computes together with every disjoint component, so the
identity $\sum_u S_u = 1$ is asserted numerically rather than assumed
(Supplement~\ref{app:sobolest}). The other quantity is retained as the
\emph{largest per-axis interaction involvement}; the two differ
substantially, at \largestInvolvementPct\ against \interactionTotalPct\ at
the median pair.

\subsection{Resolution regions and the robustness index}
\label{sec:regionsdef}

\begin{definition}[Cell labels and surface regions]
\label{def:regions}
Let $\mathcal{F}$ be the declared inference family of
Section~\ref{sec:threesurfaces} for a pair --- the cells that carry bootstrap
draws, a median \inferenceShareAdmissibleMedianPct\ of that pair's admissible
scalar cells. Under the simultaneous band of Section~\ref{sec:simbands}, a
cell of
$\mathcal{F}$ is \emph{beneficial} if its lower band is above zero,
\emph{harmful} if its upper band is below zero, and \emph{unresolved}
otherwise. A surface is \emph{uniformly beneficial} if every cell of
$\mathcal{F}$ is beneficial; \emph{conditionally beneficial} if some are and
none is harmful;
\emph{uniformly harmful} and \emph{conditionally harmful} in mirror image;
\emph{sign-changing} if there is at least one cell of each kind; and
\emph{unresolved} if none is resolved --- \textbf{or if fewer than
\minResolvedSharePct\ of $\mathcal{F}$ is resolved at all}. The
\emph{robustness index} is
\[
\rho \;=\; \frac{\#\{\text{beneficial}\} - \#\{\text{harmful}\}}
{|\mathcal{F}|} \;\in\; [-1, 1],
\]
and every directional label is reported with the counts of beneficial,
harmful and unresolved cells it rests on.
\end{definition}

\noindent \textbf{The minimum resolved share is not free, and the family the
label is relative to is not neutral.} Without the share a surface earns a
direction from as few as two resolved cells of a hundred and eighty, which on
a family the data barely touch is a statement about the band's width and not
about the register; \nLabelsLostToMinShare\ of the \nDirectionalLabels\ pairs
that carry a direction lose it under the condition (Table~\ref{tab:triple}),
and $\rho$ is printed beside the triple everywhere because $\rho$ alone cannot
distinguish two resolved cells from ninety. The threshold is a declared
convention: at \minShareLowPct\ it withdraws \nLabelsLostMinOne\ directions,
at the declared \minResolvedSharePct\ \nLabelsLostToMinShare, at
\minShareHighPct\ \nLabelsLostMinTen, and Table~\ref{tab:triple}'s source
file carries all three. And a label computed on a sixth of a pair's cells is
systematically cleaner than one computed on all of them, with \emph{uniformly
beneficial} the label most helped, so the manuscript never says a surface is
uniformly beneficial without this restriction being in force, and
Section~\ref{sec:regions} reports the one pair that reaches it together with
the share of its admissible cells the label was earned over.

\noindent The two harmful labels are not decoration, and what this corpus
does with them is a case of the minimum share doing its work.
Table~\ref{tab:triple}, which prints the label before and after the minimum
share, carries \nCondHarmful\ conditionally harmful and
\nUniformlyHarmful\ uniformly harmful surfaces out of \nPairs\ before it;
each conditionally harmful label rests on at most \nCondHarmfulResolvedMax\
resolved cells of \nHelpdeskFamily, so Definition~\ref{def:regions} withdraws
them, and under the definition this paper reports the count of either kind
is \nCondHarmfulMinShare.
The label is available and the data do not earn it. $\rho = 1$ is the only state in
which one positive number is a safe summary of a surface, and the empirical
question this paper answers is how often it obtains. The region is defined
over the \emph{scalar} family; folding the decision-curve cells into it would
let a single extreme threshold decide the label of a whole surface, so the
decision curve gets its own region (Section~\ref{sec:dca}). The
intercept-only rung is excluded from every admissible set, because a reduction
measured against a model with no features is inflated by construction; the
master surface file \emph{does} carry that rung, which is why the two counts
in Table~\ref{tab:denominator} differ.

\subsection{Specification regret, as a robustness check}
\label{sec:regretdef}

\textbf{This subsection defines a check and not a reporting object.} What
\emph{specification regret} reports is a null --- the one-number rule's median
excess regret is zero in every instrument, the out-of-sample ordering reverses
under rolling folds (Section~\ref{sec:rules}), and on calibrated net benefit
what separates the rules is the conservative one rather than the one-number
one (Section~\ref{sec:nbregret}) --- and a null is a check on
the paper's other two arguments rather than a result of its own. It is
reported in full, and it is not a contribution.

The construction asks what choosing one specification costs against the best
one, on a scale a decision can use; the identity map's maximum excess is
\identityMaxExcess. Supplement~\ref{app:regretfull} gives the three declared
designs --- metric-specific, calibrated net benefit, and partially identified
over a declared set of monotone utility maps --- and the four decision rules
compared under them.

\subsection{A reporting standard, in one paragraph}
\label{sec:standard}
Report the design space and the measure over it; the denominator the claim
quantifies over, named as one of the three; the absolute increment at every
rung rather than a ratio; the decomposition with $1 - \sum_i S_i$ reported as
such; the region and $\rho$ under the band's own family, with the counts they
rest on; and the reduction $R$ last if at all, and only where its denominator
is resolvably positive. Supplement~\ref{app:standard} states the standard in
full, together with the errors this study made before it got them right, a
six-step recipe, the \texttt{fieldvalue} package
\citep{fieldvalue} that implements the objects and refuses by construction to
emit a single number, and a worked example on a dataset that is not an event
log at all.
